%% kapa: exam=1 round=19 year=2008 period=2 form=A subject=economics subject_ko=경제원론 seq=29 number=29
%% source: Q-Net 감정평가사 기출문제 (cst00309) artlSeq=1001962 file=extracted/1008/1001962/2008년 제19회 감정평가사 1차시험 기출문제/2교시.hwp
%% conversion: kordoc 4.24.51 -> tools/gen1.py (auto); review=proofread-sonnet; flags=none
\begin{question}{29}
영희는 소득이나 화장품 가격의 변화에 관계없이 소득의 5분의 1을 화장품 값으로 지출한다. 영희의 화장품 수요에 대한 설명 중 옳은 것은?
\begin{choices}
  \item 영희에게 화장품은 사치재에 속한다.
  \item 영희에게 화장품은 열등재에 속한다.
  \item 화장품 가격이 10\% 상승하면 영희의 화장품 수요량은 10\% 감소한다.
  \item 영희의 화장품에 대한 수요곡선은 우하향하는 직선 형태를 지닌다.
  \item 영희의 소득이 10\% 증가하면 화장품 값 지출이 2\% 증가한다.
\end{choices}
\end{question}
\answer{3}
